\documentclass[a4paper,10pt]{article}

\usepackage[ngerman]{babel}
%\usepackage[top=3cm,bottom=3cm,left=3cm,right=3cm]{geometry} %NICHT IN KEMIE2


\usepackage{microtype} %schöneres typesetting  mit [stretch=10] für nur 1% anstelle von 2%
\usepackage{lmodern} %Latin Modern FONT
\usepackage{pifont} %schönere Font

\usepackage{chemfig}
\usepackage[version=4]{mhchem}
\usepackage{amsmath}
\usepackage{nicefrac}
\usepackage{amssymb}
%\usepackage{mathtools}

\usepackage{titlesec} %Um Section, etc bearbeiten zu können
\usepackage{fancyhdr}
\usepackage{setspace}
\usepackage{enumitem} %enumerate
\usepackage{arydshln} %für dashline
\usepackage{multicol}
\usepackage{lipsum}
\usepackage[labelfont=bf]{caption}

\usepackage[many]{tcolorbox}

%\usepackage{pgfornament} %Scheene Ornamente
\usepackage{witharrows} %Arrows in align-Umgebung
%\usepackage{awesomebox} %Notizboxen, etc
%\usepackage{textpos} %Für zum Beispiel vertikale Texte am Rand


%\usepackage{wrapfig} %Text um Bilder
%\usepackage{varwidth} % Alternative für engere minipage


\usepackage{hyperref}


\tcbuselibrary{theorems}

%\titlespacing*{\section}{0pt}{10mm}{3mm}
%\titlespacing*{\subsection}{0pt}{6mm}{3mm}
%\titlespacing*{\subsubsection}{0pt}{8mm}{3mm}


\renewcommand{\ttdefault}{lmtt} %Schriftart wird geändert zu Latin Modern Typewriter


\hypersetup{hidelinks,
			colorlinks=true,
			linkcolor=black,
			citecolor=miudiutex,
			urlcolor=miugray2
			}

\setlength\fboxsep{0pt}				
\setlength\columnsep{20pt}
\setlength{\parindent}{0pt}
\setcounter{section}{0}






%\definecolor{miudiutex}{HTML}{2f3d39}
\definecolor{miudiutex}{HTML}{1d2926}
\definecolor{miugray}{HTML}{b4cccf}
%\definecolor{miugray2}{HTML}{4d7365}
\definecolor{miugray2}{HTML}{6F8C81}
\definecolor{beige}{HTML}{F4F8D3}
\definecolor{white2}{HTML}{f5f7f8}
\definecolor{red}{HTML}{cf1f5c}
\definecolor{red2}{HTML}{9C191B}
\definecolor{green}{HTML}{00A36C}
\definecolor{delphinidin2}{HTML}{315794}
\definecolor{delphinidin}{HTML}{8B9EBA}
\definecolor{uniblau}{HTML}{004e9f}
\definecolor{unigelb}{HTML}{fcba00}
\definecolor{unigrau}{HTML}{909085}
\definecolor{pelargonidin}{HTML}{b33807}
\definecolor{pelargonidin2}{HTML}{FFDC96}
\definecolor{cyanidin}{HTML}{ba0746}
\definecolor{paeonidin}{HTML}{d15ca6}
\definecolor{petunidin}{HTML}{B88495}
\definecolor{petunitext}{HTML}{963354}
\definecolor{malvidin}{HTML}{8a2d8c}
\definecolor{mytheorembg}{HTML}{F2F2F9}
\definecolor{mytheoremfr}{HTML}{00007B}
\definecolor{uniblau2}{HTML}{27548A}
\definecolor{unigelb2}{HTML}{DDA853}
\definecolor{unidunkelblau2}{HTML}{183B4E}
\definecolor{unibeige}{HTML}{F5EEDC}


\tikzset{>=Latex}

\raggedcolumns %wichtig für Multicolumns
\color{miudiutex}



\usetikzlibrary{chains,
                positioning,
                shapes.multipart,
                }


\setchemfig
{%	
	bond offset=2pt,
	atom sep=15pt, 
	cram width=3pt,
	cram dash width=0.8pt,
	cram dash sep=1.2pt,
	double bond sep=2.5pt, 
	bond style={line width=1.1pt,cap=round},
	baseline=(current bounding box.center),
	bond join=true
}
\setcharge{extra sep=3pt}





%================================
% SYMBOLS
%================================

\newcommand{\RA}{$\rightarrow$~}
\newcommand{\RL}{$\rightleftharpoons$~}
\newcommand{\HARP}{\ding{226}~}
%$\xrightarrow{2}$ %Pfeil mit 2 drüber


%================================
% SHORT COMMANDS (BOXES ; ETC)
%================================


\newcommand{\info}[1]{\begin{note}{\color{miudiutex}#1}\end{note}}
\newcommand{\qsbig}[1]{\begin{qstion}{}{}\begin{center}
			#1 \end{center}\end{qstion}}
\newcommand{\notiz}[1]{\begin{notizz}{}{}\begin{center}
			{\color{miugray2!50!miudiutex}#1} \end{center}\end{notizz}}
\newcommand{\klausur}[1]{\begin{klausurinfo}{}{}\begin{center}
			#1 \end{center}\end{klausurinfo}}
\newcommand{\dfn}[1]{\begin{defin}{}{}\begin{center}
			#1 \end{center}\end{defin}}
\newcommand{\gesetz}[2]{\begin{gestex}{}{}\begin{center}
			#1 \end{center}\end{gestex}}
\newcommand{\gergesetz}[1]{\begin{gesgex}{}{}\begin{center}
			#1 \end{center}\end{gesgex}}
\newcommand{\klaausur}[1]{\begin{klausurinfo}[width=.77\linewidth]{}{}\begin{center}
			\footnotesize #1 \end{center}\end{klausurinfo}}
\newcommand{\winfo}[1]{\begin{wichtig}{}{} \begin{center}
#1 \end{center} \vspace*{0mm} \end{wichtig}}
\newcommand{\zit}[2]{\begin{zitat}{#1}{}\small #2 \end{zitat}}



\newcommand{\unten}{\end{center} \tcblower \begin{center}}


\newcommand*\circled[1]{\tikz[baseline=(char.base)]{
            \node[shape=circle,draw,inner sep=0.5pt,miudiutex] (char) {\tiny #1};}}
            
\newcommand{\miusquare}{{\color{miugray2!50}\scriptsize $\blacksquare$}~}
\newcommand{\miusquarp}{{\color{petunidin!50}\scriptsize $\blacktriangleright$}~}
\newcommand{\niu}{\item[\miusquare]}
\newcommand{\nip}{\item[\miusquarp]}
\newcommand{\mrk}[1]{{\color{petunitext!55!red}\textsc{#1}}}

%%%%% ABSAETZE

\newcommand{\pps}{\par \vspace*{1mm}}
\newcommand{\pp}{\par \vspace*{3mm}}
\newcommand{\ppbig}{\par \vspace*{8mm}}
\newcommand{\pphuge}{\par \vspace*{10mm}}



%================================
% HEADER & FOOTER FUER UNI
%================================
\newcommand{\miusfoot}{
\fancypagestyle{plain}{%
  \fancyhf{}%
  \renewcommand{\footrulewidth}{0.0mm}%
  \fancyfoot[L]{\hspace*{-3cm} {\color{miugray2}\textbf{\thepage}}}%
  \renewcommand{\headrulewidth}{0mm}%
} 
\pagestyle{plain}
}

\newcommand{\miushead}[2]{
\miusquare #1 - 4.Semester - ELW \hfill 
\includegraphics[width=60pt]{{F://LaTeX-Files/shiet/unibo}}
\par \vspace*{-3mm}
\hrulefill
\par \vspace*{1mm}
\par \vspace*{-4mm}
\normalsize	
	\begin{flushright}
	\textbf{\miusquare #2}\par \vspace*{1.5mm}
	\small	
	\textit{Letzte Überarbeitung:} \par 
	\footnotesize \today
	\end{flushright}
	\normalsize
\par \vspace*{-8mm}
}

\newcommand{\sidebar}{
\AddToHook{shipout/background}{
  \begin{tikzpicture}[remember picture,overlay, shift={(current page.south west)}]
   \path[fill=miugray2!15](0,0) rectangle (3.5cm,\paperheight);
%   \path[fill=miugray2!8](0,0) rectangle (3.5cm,\paperheight);
   \path[draw,miudiutex!50] (3.5,0) to (3.5,\paperheight);
%   \path[draw,miudiutex!50,dashed] (3.5,0) to (3.5,\paperheight);  
%   \node[opacity=0.8](a)at(12,10){\includegraphics[width=9cm]{F://LaTeX-Files/pngs/girl01}};
  \end{tikzpicture} 
}
}


\newcommand{\oddsidebar}{
\AddToHook{shipout/background}{
   \put(0in,-\paperheight){%
      \ifodd\value{page}\relax
           \begin{tikzpicture}[remember picture,overlay, shift={(current page.south west)}]
   \path[fill=miugray2!15](0,0) rectangle (3.5cm,\paperheight);
      \path[draw,miudiutex!50] (3.5,0) to (3.5,\paperheight);
  \end{tikzpicture}
      \else
           \begin{tikzpicture}[remember picture,overlay, shift={(current page.south east)}]
   \path[fill=miugray2!15](-3,0) rectangle (3cm,\paperheight);
%      \path[draw,miudiutex!50] (-3.5,0) to (-3.5,\paperheight);
  \end{tikzpicture}
      \fi
   }
}
}
%================================
% ENVIRONMENTS
%================================

\newenvironment{ctab}[2]{%
				\begin{spacing}{#1}
					\begin{center}
					\begin{tabular}{#2}}
					{\end{tabular}
					\end{center}
				\end{spacing}
				}


%================================
% Farbiger-Background-Box
%================================



\newcommand{\farbig}[1]{
\begin{tcolorbox}[%
				hbox,
				colback=miugray2!25,
				boxsep=-2pt,
				boxrule=0pt,
				arc=0pt,
				nobeforeafter,
				tcbox raise base,
				shrink tight,
				extrude by=1pt
				]
				{\color{miudiutex}{#1}}
\end{tcolorbox}~}

%================================
% SIMPLE SCHWARZER RAHMEN BOX
%================================


\newcommand{\boxfr}[1]{
	\begin{tcolorbox}[
	colframe=miugray2,
	colback=white!95!miugray2,
	boxrule = 0.5pt, 
%	toprule=0pt, bottomrule=0pt,rightrule=0pt,leftrule=0pt,
	boxsep=0pt,
	arc=4pt,
	drop shadow southeast,
	enhanced
	]
	\color{miudiutex}
	\begin{center}
	#1
	\end{center}
	\end{tcolorbox}
}

%================================
% Antwort-Box
%================================

\newcommand{\antwort}[1]{\par 
	\begin{tcolorbox}[
	colback=gray!15,
	breakable,
	enhanced,
	frame hidden,
	arc=5pt,
	boxrule = 2pt,
%	borderline west = {2pt}{0pt}{petunidin},
%	drop shadow=miugray2!30
	]
	\color{miudiutex}	#1
	\end{tcolorbox}
}

%================================
% Frame-IT
%================================

\newcommand{\frameit}[1]{\par 
	\begin{tcolorbox}[
	colback=gray!10,
	breakable,
	enhanced,
	frame hidden,
	arc=5pt,
	boxrule = 0pt,
	borderline west = {4pt}{0pt}{petunidin},
%	drop shadow=miugray2!30
	]
		#1
	\end{tcolorbox}
}



%================================
% Question BOX
%================================

\makeatletter
\newtcbtheorem{qstion}{Frage}{enhanced,
	breakable,
	colback=white,
	colframe=delphinidin!100,
%	capture=hbox,
	boxsep=-3pt,
	attach boxed title to top left={yshift*=-\tcboxedtitleheight},
	fonttitle=\bfseries,
%	title={#2},
	boxed title size=title,
	boxed title style={%
			sharp corners,
			rounded corners=northwest,
			colback=tcbcolframe,
			boxrule=0pt,
		},
	underlay boxed title={%
			\path[fill=tcbcolframe] (title.south west)--(title.south east)
			to[out=0, in=180] ([xshift=5mm]title.east)--
			(title.center-|frame.east)
			[rounded corners=\kvtcb@arc] |-
			(frame.north) -| cycle;
		},
	#1
}{def}
\makeatother



%================================
% SIMPLE Question BOX
%================================

\newcommand{\qssmall}{
\begin{tcolorbox}[hbox,colback=miugray2!55,boxsep=-4pt,boxrule=0pt,arc=2pt,nobeforeafter, tcbox raise base,shrink tight,extrude by=3pt]
\textbf{\color{white}{?}}
\end{tcolorbox}~~}

\newcommand{\folie}[1]{
\begin{tcolorbox}[%
				hbox,
				colback=miugray2!20,
				boxsep=2pt,
				boxrule=0.5pt,
				arc=2pt,
				nobeforeafter,
				tcbox raise base,
				shrink tight,
				extrude by=3pt]
				{\color{miudiutex}{Folie #1}}
\end{tcolorbox}~}


%================================
% NOTIZ BOX
%================================

\makeatletter
\newtcbtheorem{notizz}{Notiz}{enhanced,
	breakable,
	coltitle=white,
	colback=petunidin!10,
	colframe=petunidin!50,
	attach boxed title to top left={yshift*=-\tcboxedtitleheight},
	fonttitle=\bfseries,
%	title={#2},
	boxed title size=title,
	boxed title style={%
			sharp corners,
			rounded corners=northwest,
			colback=tcbcolframe,
			boxrule=0pt,
		},
	underlay boxed title={%
			\path[fill=tcbcolframe] (title.south west)--(title.south east)
			to[out=0, in=180] ([xshift=5mm]title.east)--
			(title.center-|frame.east)
			[rounded corners=\kvtcb@arc] |-
			(frame.north) -| cycle;
		},
		#1
}{def}
\makeatother



%================================
% Question BOX: Klausur
%================================

\makeatletter
\newtcbtheorem[no counter]{klausurinfo}{Klausurrelevant}{enhanced,
%	breakable,
	colback=miugray2!10,
	colframe=miugray2!75,
	attach boxed title to top left={yshift*=-\tcboxedtitleheight},
	fonttitle=\bfseries,
	title={#2},
	arc=4pt,
	boxrule=1pt,
	boxed title size=title,
	boxed title style={%
			sharp corners,
			rounded corners=northwest,
			colback=tcbcolframe,
			boxrule=0pt,
			arc=4pt,
		},
	underlay boxed title={%
			\path[fill=tcbcolframe] (title.south west)--(title.south east)
			to[out=0, in=180] ([xshift=5mm]title.east)--
			(title.center-|frame.east)
			[rounded corners=\kvtcb@arc] |-
			(frame.north) -| cycle;
		},
	#1
}{def}
\makeatother


%================================
% DEFINITION
%================================

\makeatletter
\newtcbtheorem[no counter]{defin}{Definition}{enhanced,
	breakable,
	colback=gray!10,
	colframe=gray,
	attach boxed title to top left={yshift*=-\tcboxedtitleheight},
	fonttitle=\bfseries,
	title={#2},
	arc=4pt,
	boxrule=1pt,
	fontlower=\footnotesize,
	collower=miugray2!50!miudiutex,
	boxed title size=title,
	boxed title style={%
			sharp corners,
			rounded corners=northwest,
			colback=tcbcolframe,
			boxrule=0pt,
			arc=4pt,
		},
	underlay boxed title={%
			\path[fill=tcbcolframe] (title.south west)--(title.south east)
			to[out=0, in=180] ([xshift=5mm]title.east)--
			(title.center-|frame.east)
			[rounded corners=\kvtcb@arc] |-
			(frame.north) -| cycle;
		},
	#1
}{def}
\makeatother

%================================
% Gesetzes-Text EU
%================================

\makeatletter
\newtcbtheorem[no counter]{gestex}{EU - Verordnung}{enhanced,
	breakable,
	colback=gray!10,
	colframe={gray!65!delphinidin},
	attach boxed title to top left={yshift*=-\tcboxedtitleheight},
	fonttitle=\bfseries,
	arc=4pt,
	boxsep=10pt,
	boxrule=1pt,
	fontlower=\footnotesize,
	collower=miugray2!50!miudiutex,
	boxed title size=title,
	overlay={%
			\node[xshift=-1cm,yshift=-1mm] at (frame.north east)
			{\includegraphics[width=8mm]{F://LaTeX-Files/shiet/eu}};
%			\node[draw,miugray2] at (frame.north) {{#2}}; 
			},%
	boxed title style={%
			sharp corners,
			rounded corners=northwest,
			colback=tcbcolframe,
			boxrule=0pt,
			arc=4pt,
		},
	underlay boxed title={%
			\path[fill=tcbcolframe] (title.south west)--(title.south east)
			to[out=0, in=180] ([xshift=5mm]title.east)--
			(title.center-|frame.east)
			[rounded corners=\kvtcb@arc] |-
			(frame.north) -| cycle;
		},
	#1
}{def}
\makeatother


%================================
% Gesetzes-Text Ger
%================================

\makeatletter
\newtcbtheorem[no counter]{gesgex}{Gesetzestext}{enhanced,
	breakable,
	colback=gray!10,
	colframe={gray!95!pelargonidin},
	attach boxed title to top left={yshift*=-\tcboxedtitleheight},
	fonttitle=\bfseries,
	title={#2},
	arc=4pt,
	boxrule=1pt,
	fontlower=\footnotesize,
	collower=miugray2!50!miudiutex,
	boxed title size=title,
	overlay={\node[xshift=-1cm,yshift=-3mm] at (frame.north east)
			{\includegraphics[width=8mm]{F://LaTeX-Files/shiet/ger}}; 
			},
	boxed title style={%
			sharp corners,
			rounded corners=northwest,
			colback=tcbcolframe,
			boxrule=0pt,
			arc=4pt,
		},
	underlay boxed title={%
			\path[fill=tcbcolframe] (title.south west)--(title.south east)
			to[out=0, in=180] ([xshift=5mm]title.east)--
			(title.center-|frame.east)
			[rounded corners=\kvtcb@arc] |-
			(frame.north) -| cycle;
		},
	#1
}{def}
\makeatother


%================================
% Question BOX: Wichtig
%================================

%\makeatletter
%\newtcbtheorem[no counter]{wichtig}{Wichtige Info}{enhanced,
%	breakable,
%	colback=pelargonidin2!80,
%	colframe=red2!100,
%	fonttitle=\bfseries,
%	title={#2},
%	leftrule=5mm,
%	boxed title size=title,
%	boxed title style={%
%			sharp corners,
%			rounded corners=northwest,
%			colback=tcbcolframe,
%			boxrule=0pt,
%		},
%	underlay boxed title={%
%			\path[fill=tcbcolframe] (title.south west)--(title.south east)
%			to[out=0, in=180] ([yshift=1mm,xshift=7mm]title.east)--
%			([yshift=1mm]title.center-|frame.east)
%			[rounded corners=\kvtcb@arc] |-
%			(frame.north) -| cycle;
%		},
%		watermark tikz={\draw[line width=2mm] circle (1.1cm)
%		node{\fontfamily{ptm}\fontseries{b}\fontsize{20mm}{20mm}\selectfont !};}],
%	attach boxed title to top left={xshift=5mm,yshift*=-\tcboxedtitleheight},
%	watermark zoom=0.65,
%	  underlay={%
%    \path[draw=none] (interior.south west) rectangle node[white]{\Huge \textbf{!}} ([xshift=-4mm]interior.north west);
%    },
%	#1
%}{def}
%\makeatother


\makeatletter
\newtcbtheorem[no counter]{wichtig}{Wichtige Info}{enhanced,
	breakable,
	colback=pelargonidin2!30,
	colframe=petunidin!100,
	fonttitle=\footnotesize\bfseries,
	title={#2},
	leftrule=4mm,
	boxed title size=title,
	boxed title style={%
			sharp corners,
			rounded corners=northeast,
			colback=tcbcolframe,
			boxrule=0pt,
		},
	underlay boxed title={%
			\path[fill=tcbcolframe] (title.south east)--(title.south west)
			to[out=180, in=0] ([yshift=2mm,xshift=-7mm]title.west)--
			([yshift=2mm]title.center-|frame.west)
			[rounded corners=\kvtcb@arc] |-
			(frame.north) -| cycle;
		},
		watermark tikz={\draw[line width=2mm] circle (1.1cm)
		node{\fontfamily{ptm}\fontseries{b}\fontsize{20mm}{20mm}\selectfont !};},
	attach boxed title to top right={yshift*=-\tcboxedtitleheight},
	watermark zoom=1.45,
	watermark opacity=0.35,
	  underlay={%
    \path[draw=none] (interior.south west) rectangle node[white]{\LARGE \textbf{!}} ([xshift=-4.3mm]interior.north west);
    },
    #1
}{def}
\makeatother





%================================
% ZITAT BOX
%================================

\tcbuselibrary{theorems,skins,hooks}
\newtcbtheorem{zitat}{Zitat}
{%
	enhanced,
	breakable,
	colback = gray!10!white,
	frame hidden,
	boxrule = 0sp,
	borderline west = {2pt}{0pt}{petunidin},
	sharp corners,
	detach title,
	before upper = \tcbtitle\par\smallskip,
	coltitle = gray!50!black,
	fonttitle = \bfseries\sffamily,
	description font = \mdseries,
	separator sign none,
	segmentation style={solid, gray!50!black},
}
{th}



\input{F://LaTeX-Files/shiet/vl.tex}
\usepackage{pgfplots}
\usepackage{awesomebox}
%\usepackage{sidenotes}
%\setcounter{section}{1} % <--- FALLS KEINE SECTION DEFINIERT WIRD2

\begin{document}

\sidebar
\miushead{Vorlesung 13 + 14 - PLMT}{Vorlesung 13 + 14}
\par \vspace*{5mm}
\section{Vorlesung 13 - Verdampfen}

\subsection{Phasendiagramm}
Um Wasser aus der flüssigen Phase in die gasförmige Phase zu überführen, können wir dies durch eine \textbf{Temperaturerhöhung} oder \textbf{Druckerniedrigung} realisieren.



%\pp
%Wie überführen wir flüssiges Wasser in die gasförmige Phase?
%\begin{itemize}
%\item Durch Temperaturerhöhung
%\item Durch Druckerniedrigung
%\end{itemize}


%\begin{marginfigure}
\begin{figure}[h]
\centering
\begin{tikzpicture}[scale=0.8]
\draw[->] (0,0) to (4,0) node[right] {T};
\draw[->] (0,0) to (0,4) node[right] {P};
\draw[-, rounded corners=5pt] (0.5,0.5) to (1,1) to (1,3);
\draw[-, petunitext] (1,1) to [bend right=35](3,3);
\draw[->, miugray2] (1.5,2) to (3.5,2);
\draw[->, miugray2] (1.5,2) to (1.5,0.5);
\node[scale=0.7] at(1.75,3) {flüssig};
\node[scale=0.7] at(3,1) {gasförmig};
\end{tikzpicture}
\caption{Dampfdruckkurve}
\end{figure}
%\end{marginfigure}
\footnotesize
\textit{Ein Phasendiagramm (auch Zustandsdiagramm, Zustandsschaubild oder Gleichgewichtsschaubild) ist eine Projektion von Phasengrenzlinien aus dem Zustandsraum eines thermodynamischen Systems in ein zweidimensionales kartesisches Koordinatensystem. Verschiedene Zustände lassen sich über die Clausius-Clapeyron-Gleichung berechnen.}
\normalsize



\subsection{Energie für Verdampfung berechnen}

Für die Verdampfung benötigte Wärmezufuhr $\dot{Q}_H$:
\pp
\frameit{
\color{petunitext}
\begin{align}
\dot{Q}_H = \dot{m}_D \cdot \Delta h_v + \dot{m}_F \cdot c_{p,F} \cdot (T_S - T_F)
\end{align}
\color{miudiutex}
}
\reversemarginpar
\marginnote[\small Der Begriff \textbf{Brüden} bezeichnet in der Regel mit Wasserdampf gesättigte Luft, die beim Trocknen von Feststoffen entsteht.]{}
\pp
Wobei: \par
\begin{multicols}{2}
$\dot{m}_D$ = Massenstrom des Brüden \par
$\dot{m}_F$ = Massenstrom des Zulaufs (F \RA feed) \par 
$T_S$ = Siedetemperatur \par 
$T_F$ = Temperatur des Zulaufs \par
\end{multicols}
\pagebreak
\subsection{Sattdampf}
\begin{itemize}
\item $T$, $p$ durch Dampfdruckkurve definiert
\item keine Nebeltröpfchen
\item Wärmeentzug \RA Bildung von Nebeltröpfchen durch Kondensation \RA \textbf{Nassdampf}
\item Wenn Wärmeentzug sehr vorsichtig vollzogen wird \RA \textbf{übersättigter Dampf}
\item Bei erhitzen \RA \textbf{Überhitzter Dampf} (Brüden)
\end{itemize}

\subsection{Dampfdruckkurve}
Beschreibung der Dampfdruckkurve durch die Beziehung von \textbf{Clausius-Clapeyron}:
\begin{align*}
\ln \frac{p_2}{p_1} = \frac{\Delta h_v M_L}{R} \cdot ( \frac{1}{T_1} - \frac{1}{T_2} )
\end{align*}

Wobei: \par
\begin{multicols}{2}
$p$ = Dampfdruck \par 
$R$ = Universelle Gaskonstante \par 
$\Delta h_v$ = mittlere spezifische Verdampfungsenthalpie \par 
$M_L$ = molare Masse des Lösungsmittels \par 
$\Delta h_{vM}$ = mittlere molspezifische Verdampfungsenthalpie
\end{multicols}
\pp
Nach Umformen erhalten wir
\begin{align*}
\Delta h_{vM} = \frac{R \cdot T_1 \cdot T_2 \cdot \ln(\frac{p_2}{p_1})}{T_2 - T_1}
\end{align*}

Die oben stehende Formel lässt sich unter der Annahme, dass ein ideales Gas vorliegt und das Molvolumen eines Gases größer als das einer Flüssigkeit ist, vereinfachen und es gilt annäherungsweise\par 
\frameit{
\color{petunitext}
\begin{align}
\Delta h_{vM} = \frac{R \cdot T^2}{\Delta T} \cdot \frac{\Delta p}{p_0}
\end{align}
\color{miudiutex}
}

\pagebreak
\subsection{Siedepunkterhöhung}
\subsubsection{Durch gelösten Stoff} 
\textit{Beispiel: Salz in Nudelwasser} \par 
Wir können eine Siedepunkterhöhung durch das Lösen eines Stoffes (mit vernachlässigbarem kleinem Dampfdruck) erreichen. \par Die Dampfdruckerniedrigung $\Delta p$ ist proportional zum Stoffmengenanteil $x_i$ des gelösten Stoffes $i$
\reversemarginpar
\marginnote[\small Raoultsches Gesetz zur Berechnung des Gas-Flüssig-Gleichgewichts \RA]{}[1cm]
\pp
\frameit{
\color{petunitext}

\begin{align}
\Delta p = x_i \cdot p
\end{align} 
}
\par \vspace*{2mm}
\color{miudiutex}

Der Stoffmengenanteil $x_i$ berechnet sich aus dem Verhältnis von unserem Stoff $i$ im Verhältnis zur Gesamten Stoffmenge:
\begin{align*}
x_i = \frac{n_i}{n_{tot}}
\end{align*}


Für Salze gilt:
\begin{align*}
n_i = n_{Salz} \cdot [1+ \delta \cdot (z-1)]
\end{align*}
Wobei: \par 
$\delta$ = Dissoziationsgrad des Salzes \par 
$z$ = Anzahl Ionen pro Formeleinheit des Salzes \pp

\ppbig
Formel zur Berechnung der Temperaturdifferenz des Siedepunkts bei Zugabe eines Lösungsmittels:

\begin{align*}
\Delta T &= \frac{R \cdot T^2}{\Delta h_{vM]}} \cdot \frac{\Delta p}{p_0} \\
\Delta T &= \frac{R \cdot T^2}{\Delta h_{vM]}} \cdot x_i \\
&\approx \frac{R \cdot T^2}{\delta_L \cdot \Delta h_v} \cdot c_i
\end{align*}
\frameit{
\color{petunitext}
\begin{align}
\Delta T &\approx \frac{R \cdot T^2}{\delta_L \cdot \Delta h_v} \cdot c_i
\end{align}
\color{miudiutex}
}
\pp
Wobei:\par 
$\delta_L$ = Dichte des Lösungsmittels \par 
$c_i$ = molare Konzentration des gelösten Stoffes i [mol/m$^3$]

\ppbig 


\subsubsection{Durch hydrostatischen Druck}
\textit{Beispiel: Kochen im Vakuum oder Druckkessel}
\begin{align*}
\Delta T = \frac{R \cdot T^2}{\Delta h_{vM}} \cdot \frac{\rho_{lg} \cdot g \cdot h}{p}
\end{align*}
Wobei: \par
$\rho_{lg}$ = Dichte der Flüssigkeits-Dampf-Säule mit Dampfanteil $\varepsilon$
\begin{align*}
\rho_{lg} = (1 - \varepsilon) \cdot \rho_l
\end{align*}
$\rho_l$ = Dichte der Flüssigkeit ist






\ppbig





\subsection{Beispielvideos}
Thin Film Evaporation Process \par 
LCI Corporation \pp
Pfaudler Wiped Film Evaporator (WFE) - Construction and Operation \par 
GMM Pfaudler \pp
Falling Film Evaporator \par 
gea.com \pp







\subsection{Brüdenverdichtung}
\begin{figure}[h]
\centering

\begin{tikzpicture}
\draw (0,0) to (0,1);
\draw (-1,0.8) to (0,0.8);
\draw (0,2.5) to (0,1.5) to (0.25,1) to (0.25,0.8);
\draw[xscale=-1,xshift=-1cm] (0,2.5) to (0,1.5) to (0.25,1)to (0.25,0.8);
\draw (1,0) to (1,1);
\draw[->,miugray2] (0.5,2.5) to (0.5,1.5) node[right,xshift=5mm] {$T_1$, $p_1$};
\draw[->,petunitext] (-1.3,1.1) to (-0.6,1.1) node[above,xshift=-3mm,scale=0.8,yshift=1mm] {Brüden};
\draw[dashed,rounded corners=5pt,->,petunitext!35] (-0.5,1.1) to (0.1,1.1) to (0.25,0.3);
\draw (-1,1.6) to (0,1.6);
\end{tikzpicture}
\caption{Brüdenverdichtung}
\label{fig1}
\end{figure}

In \textbf{Abbildung \ref{fig1}} wird gezeigt, wie durch eine Verengung eines Rohres Wasserdampf fließt und ein Unterdruck entsteht. Der Brüden der von der Seite hinzuströmt wird ''eingesogen'' und durch diese Kraft verdichtet.\par
Durch die Verdichtung des Brüden lässt sich ein hoher Teil der Energie, die ursprünglich zum Erhitzen des Lebensmittels aufgewendet wurde, zurückgewinnen. Eine Folie aus der Vorlesung zeigt den Mechanismus der Energierückgewinnung in einem Mehrstufenprinzip.

\pagebreak
\section{Vorlesung 14 - Trocknen}
\subsection{Welche Lebensmittel trocknen wir?}
An der Tafel wurden sämtliche Ideen der Studenten gesammelt: \par 
Chips, Obst, Gemüse, Gewürze, Fleisch, Fisch, Kaffee, Tee, Kakao, Backwaren, Teigwaren, Milch, Getreide, Algen, Hülsenfrüchte, etc.
\pp
\boxfr{
Das Trocknen bzw. die Haltbarmachung durch Wasserentzug ist eines der ältesten Haltbarmachungsmethoden der Menschheit. Welche Gründe des Trocknens kennen wir noch?
}

\subsection{Effekte des Trocknens}
\subsubsection{Erwünschte Effekte}
\begin{itemize}
\item Konservierung, Haltbarmachung
\item Reduktion von Gewicht und Volumen
\item Änderung der Struktur und technofunktionellen Eigenschafte
\end{itemize}

\subsubsection{Unerwünschte Effekte}
\begin{itemize}
\item Verlust bzw. Veränderung von wertgebenden Inhaltsstoffen / Eigenschaften
\begin{itemize}
\item Vitamine
\item Textur
\item Aroma
\item Farbe
\item Volumen, Form
\end{itemize}
\end{itemize}

\subsection{Abtrennen von Wasser}

\begin{multicols}{2}

\subsubsection{Aus feuchter Luft}
\begin{itemize}
\item Adsorption (zum Beispiel wenn ein Lebensmittel Wasser aus der Luft aufnimmt)
\item Kondensation (Spiegel im Badezimmer)
\item Bildung von Kristallwasser (CaCl$_2$)
\end{itemize}

\columnbreak
\subsubsection{Aus wasserhaltigen Flüssigkeiten}
\begin{itemize}
\item Verdampfen (Destillation)
\item Zentrifugieren
\item Bildung von Kristallwasser
\end{itemize}
\end{multicols}

\subsubsection{Mechanische Verfahren}
\begin{itemize}
\item thermische Verfahren \RA Verdampfung oder Sublimationen
\end{itemize}

\pagebreak
\subsection{Beispiel}
Wir wollen einen Apfelring trocknen. Wie gehen wir vor?

\begin{figure}[h]
\centering
\begin{tikzpicture}[node distance=1.5cm]
\node (waermezufuhr) {Wärmezufuhr};
\node (temperaturerhoehung) [below of=waermezufuhr] {Temperaturerhöhung};
\node (dampfdruckerhoehung) [below of=temperaturerhoehung] {Dampfdruckerhöhung};
\node (verdampfung) [below of=dampfdruckerhoehung] {Verdampfung};
\node (temperaturabsenkung) [below of=verdampfung] {Temperaturabsenkung};

\draw[->] (waermezufuhr) -- (temperaturerhoehung);
\draw[->] (temperaturerhoehung) -- (dampfdruckerhoehung);
\draw[->] (dampfdruckerhoehung) -- (verdampfung);
\draw[->] (verdampfung) -- (temperaturabsenkung);

\node (1) [below of=temperaturabsenkung,petunitext, align=center, scale=0.7,yshift=1.5cm] {\RA begünstigt \\ Wärmetransport};
\end{tikzpicture}
\caption{Abfolgende Schritte eines Trocknungsprozess}
\end{figure}

Der Erfolg unseres Trockungsprozesses des Apfelrings hängt vorallendingen von Faktoren wie dem Verhältnis zur Oberfläche (zum Volumen), Temperatur, Druck, relative Feuchte und der Zusammensetzung des Produktes ab.







\end{document}